% --- BẮT ĐẦU PHẦN TỔNG KẾT ---

\chapter{Tổng kết}

Dự án ``Hệ thống Quản lý Thư viện Trực tuyến tích hợp AI'' được xây dựng nhằm
số hóa các nghiệp vụ thư viện, đồng thời cải thiện trải nghiệm tra cứu và gợi ý
sách cho người dùng. Sau quá trình phân tích, thiết kế, hiện thực và kiểm thử,
nhóm đã hoàn thiện một hệ thống có đầy đủ ba lớp chính: Back-end nghiệp vụ,
Front-end người dùng và AI Service độc lập.

\section{Kết quả đạt được}

\subsection{Back-end nghiệp vụ}

Back-end được xây dựng bằng Spring Boot theo hướng \textbf{Modular Monolith},
tách thành các module có ranh giới rõ ràng:

\begin{itemize}
    \item \textbf{Auth/User:} Đăng ký, đăng nhập, JWT, refresh token, xác thực email,
    OAuth2, hồ sơ người dùng, phân quyền theo vai trò.
    \item \textbf{Catalog:} Quản lý ấn phẩm, tác giả, thể loại, nhà xuất bản, bản sao
    sách, upload tài liệu và metadata phục vụ AI.
    \item \textbf{Circulation:} Mượn/trả sách, đặt trước, xác nhận nhận sách, phí phạt,
    thanh toán, dashboard thủ thư và các tác vụ định kỳ.
    \item \textbf{Recommendation:} Wishlist, đánh giá sách, lịch sử tìm kiếm, gọi AI
    Gateway, nhận callback xử lý AI có kiểm tra chữ ký.
    \item \textbf{Shared Infrastructure:} Exception handling, audit log, email, storage,
    Kafka/event, security annotation và các port dùng chung.
\end{itemize}

Cách tổ chức này giúp hệ thống giữ được cấu trúc rõ ràng, dễ mở rộng, nhưng
không tạo thêm độ phức tạp vận hành như mô hình microservices hoàn chỉnh.

\subsection{Front-end người dùng}

Front-end được hiện thực bằng \textbf{React, TypeScript và Vite}. Ứng dụng chia
rõ các vùng chức năng:

\begin{itemize}
    \item \textbf{Trang công khai:} Trang chủ, tìm kiếm, chi tiết sách, đăng nhập,
    đăng ký, hướng dẫn và các trang thông tin.
    \item \textbf{Khu vực độc giả:} Dashboard, sách đang mượn, đặt trước, phí phạt,
    wishlist, hồ sơ, thông báo và tìm kiếm.
    \item \textbf{Khu vực thủ thư/admin:} Quản lý đầu sách, bản sao, giao dịch,
    lưu thông, người dùng, chính sách và dashboard vận hành.
    \item \textbf{Trải nghiệm hệ thống:} Tự động refresh token, đa ngôn ngữ, dark mode,
    thông báo realtime, upload panel và tích hợp các API AI.
\end{itemize}

\subsection{AI Service}

AI Service được tách thành một dịch vụ độc lập bằng \textbf{Python FastAPI} và
worker nền. Các chức năng chính gồm:

\begin{itemize}
    \item \textbf{Xử lý PDF:} Trích xuất văn bản bằng PyMuPDF, làm sạch nội dung,
    chia chunk và bỏ qua xử lý lại khi file không thay đổi.
    \item \textbf{Tìm kiếm ngữ nghĩa:} Sinh embedding bằng SBERT tiếng Việt, lưu vector
    vào PostgreSQL/pgvector và truy vấn nearest-neighbor cho semantic search.
    \item \textbf{Metadata AI:} Dùng Gemini ở bước reduce để tạo tóm tắt, đối tượng
    độc giả và tag; có fallback deterministic khi LLM lỗi hoặc hết quota.
    \item \textbf{Gợi ý sách:} Kết hợp ALS implicit feedback, content-based realtime
    và trending fallback để phục vụ cả người dùng mới lẫn người dùng có lịch sử.
    \item \textbf{Tích hợp bất đồng bộ:} Worker nhận message qua RabbitMQ/Celery và
    callback kết quả về Back-end.
\end{itemize}

\subsection{Kiểm thử và chất lượng}

Nhóm đã xây dựng bộ kiểm thử ở nhiều lớp:

\begin{itemize}
    \item Back-end có 79 test case trên 23 test class, bao phủ auth, catalog,
    circulation, recommendation và các luồng tích hợp với PostgreSQL qua Testcontainers.
    \item Front-end có Jest và Testing Library để kiểm thử service contract, xử lý lỗi,
    ngôn ngữ, theme và các luồng giao diện quan trọng.
    \item AI Service có pytest cho PDF processing, LLM client, SQL contract, worker
    contract, recommender engine và API contract.
\end{itemize}

\subsection{Đóng gói và hạ tầng}

Các thành phần chính đã được chuẩn bị Dockerfile và Docker Compose để phục vụ
chạy thử, kiểm thử và định hướng triển khai. Nhóm cũng đã nghiên cứu mô hình
triển khai với NGINX Reverse Proxy, biến môi trường, DNS, HTTPS/SSL và tách biệt
môi trường build/runtime bằng multi-stage build.

\section{Đánh giá}

Nhìn chung, hệ thống đã đạt được các mục tiêu quan trọng của đồ án:

\begin{itemize}
    \item Hoàn thiện lõi nghiệp vụ thư viện từ quản lý sách đến mượn/trả, đặt trước,
    phí phạt và thông báo.
    \item Tích hợp AI theo hướng thực dụng, có API riêng, xử lý bất đồng bộ và có
    cơ chế fallback khi dịch vụ AI bên ngoài không ổn định.
    \item Front-end có đủ giao diện cho độc giả, thủ thư và admin, đồng thời kết nối
    được với các luồng nghiệp vụ chính.
    \item Kiến trúc tổng thể đủ rõ để tiếp tục mở rộng mà không cần viết lại toàn bộ
    hệ thống.
\end{itemize}

\section{Hạn chế}

Một số điểm vẫn cần tiếp tục hoàn thiện:

\begin{itemize}
    \item Chất lượng gợi ý phụ thuộc vào lượng tương tác của người dùng, nên vẫn còn
    bài toán cold-start cho tài khoản hoặc đầu sách mới.
    \item Pipeline AI phụ thuộc vào quota và độ ổn định của Gemini, dù đã có fallback
    để giảm rủi ro.
    \item Dữ liệu kiểm thử hiện còn giới hạn, chưa phản ánh đầy đủ quy mô vận hành của
    một thư viện lớn.
\end{itemize}

\section{Hướng phát triển}

Trong các phiên bản tiếp theo, hệ thống có thể được mở rộng theo các hướng:

\begin{itemize}
    \item Bổ sung cơ chế incremental indexing để cập nhật vector nhanh hơn khi sách
    hoặc tài liệu thay đổi.
    \item Cải thiện mô hình gợi ý bằng thêm tín hiệu hành vi và đánh giá thực tế từ
    người dùng.
    \item Mở rộng dashboard thống kê cho thủ thư và admin, hỗ trợ ra quyết định về
    tồn kho, nhu cầu mượn và chất lượng tài liệu.
\end{itemize}

\section{Kết luận}

Dự án đã xây dựng được một hệ thống quản lý thư viện trực tuyến có cấu trúc tốt,
có khả năng xử lý các nghiệp vụ cốt lõi và tích hợp AI ở mức có thể vận hành thử
nghiệm. Kết quả này tạo nền tảng vững chắc để nhóm tiếp tục hoàn thiện triển
khai thực tế, tối ưu chất lượng gợi ý và mở rộng hệ thống cho các nhu cầu thư
viện quy mô lớn hơn.
